% =====================================================================
%  Fig. 1 - Fluxo da analise (TikZ). Incluido em main.tex com \input.
%  Edite os textos diretamente aqui; as cores estao definidas em main.tex
%  (corEscrita, corLeitura).
% =====================================================================
\begin{tikzpicture}[
    font=\footnotesize,
    >={Stealth[length=2mm]},
    caixa/.style={draw=black!55, rounded corners=2pt, align=center,
                  inner sep=3pt, line width=0.5pt},
    base/.style={caixa, text width=2.9cm, minimum height=1.25cm},
    meio/.style={caixa, text width=3.1cm, fill=black!4},
    metodo/.style={caixa, text width=4.45cm, minimum height=1.25cm,
                   fill=white, align=left},
    seta/.style={->, line width=0.6pt, draw=black!70},
    rotulo/.style={font=\scriptsize\bfseries, text=black!65}
]
% ---------------- bases de dados
\node[base, fill=corEscrita!10, draw=corEscrita] (keyrecs)
  {\textbf{Escrita · KeyRecs}\\ 96 participantes\\ alvo: hesitação, em s};
\node[base, fill=corLeitura!10, draw=corLeitura, below=0.35cm of keyrecs] (complex)
  {\textbf{Leitura · CompLex}\\ 8.579 termos em contexto\\ alvo: termo difícil};

% ---------------- modelos preditivos
\node[meio, right=0.75cm of $(keyrecs.east)!0.5!(complex.east)$, anchor=west,
      minimum height=2.85cm] (ativ)
  {\textbf{Modelo preditivo}\\[2pt]
   Interpretável por construção:\\ regressão, árvore, EBM\\[1pt]
   $\times$\\[1pt]
   Caixa-preta:\\ Random Forest, XGBoost\\[2pt]
   validação cruzada $5\times10$};

% ---------------- 4 perguntas, 7 metodos globais
\node[metodo, right=0.8cm of ativ.north east, anchor=north west, yshift=0.0cm] (m12)
  {\textbf{Como as variáveis se combinam?}\\
   \textcircled{\scriptsize 1} H de Friedman\\
   \textcircled{\scriptsize 2} Decomposição funcional};
\node[metodo, right=0.2cm of m12] (m34)
  {\textbf{Quais variáveis importam?}\\
   \textcircled{\scriptsize 3} Permutação\\
   \textcircled{\scriptsize 4} Remoção, LOFO};
\node[metodo, below=0.35cm of m12] (m56)
  {\textbf{Como resumir o forest?}\\
   \textcircled{\scriptsize 5} Substitutos globais\\
   \textcircled{\scriptsize 6} Extração de regras};
\node[metodo, below=0.35cm of m34] (m7)
  {\textbf{Quais casos são atípicos?}\\
   \textcircled{\scriptsize 7} Protótipos e críticas,\\
   \phantom{x}MMD-critic};

\node[draw=black!40, dashed, rounded corners=3pt, inner sep=4pt,
      fit=(m12)(m34)(m56)(m7)] (a4) {};
\node[rotulo, anchor=south west] at (a4.north west)
  {MÉTODOS GLOBAIS APLICADOS AO RANDOM FOREST DE CADA TAREFA};

% ---------------- setas
\draw[seta] (keyrecs.east) -- ++(0.35,0) |- (ativ.west);
\draw[seta] (complex.east) -- ++(0.35,0) |- (ativ.west);
\draw[seta] (ativ.east) -- (a4.west |- ativ.east);
\end{tikzpicture}
